\clearpage\fancyhead[L]{\small PALLAS.}\hypertarget{native-374-lower}{}\label{native-374-lower}
\begin{center}\textbf{PALLAS.}\par\vspace{4pt}\hrule\vspace{4pt}\fontsize{10}{12}\selectfont Göttingische gelehrte Anzeigen. Stück 176. Seite 1753\;..\;1755. 1813. November 4.\par\vspace{4pt}\hrule\vspace{4pt}\end{center}

\par\indent Die schöne Bestätigung, welche die auf die Störungsrechnungen gegründete Vorausbestimmung
der Bewegung der \emph{Pallas} durch die Beobachtungen des vorigen Jahres erhielt, und von der
wir damals im 67. und 127. Stück dieser Blätter [S. 349 u. 356 d. B.] Rechenschaft gegeben
haben, erhöhete das Interesse, womit die Astronomen die Wiedererscheinung dieses Planeten im
gegenwärtigen Jahre erwarteten. Auf der hiesigen Sternwarte wurde der Planet zum ersten Male
den 28. Juni wieder beobachtet, und genau auf dem Platze der schon zwei Jahre zuvor durch
Hrn. \textsc{Nicolai} berechneten Ephemeride gefunden. Um die Zeit der Opposition. suchte
Prof. \textsc{Gauss} so viele Beobachtungen, als nur möglich war, zu erhalten: allein das im
verwichenen Sommer so ausgezeichnet ungünstige Wetter erlaubte nur eine kärgliche Ausbeute.
Inzwischen sind diese Beobachtungen, wenn gleich bei dem schwachen Lichte des Planeten ---
auch diesesmal nur von der 10. Grösse --- nur am Kreismicrometer angestellt, vorzüglich gut
ausgefallen. Die sämmtlichen Beobachtungen des Prof. \textsc{Gauss} sind folgende:

\begin{center}\fontsize{10}{12}\selectfont\setlength{\tabcolsep}{5pt}\renewcommand{\arraystretch}{1.15}\begin{tabular}{l|c|c|c}
1813 & Mittlere Zeit & Gerade Aufsteig. & Nordl. Abweich.\\\hline
Juni 28 & $11^{\mathrm u}\;37^{\mathrm m}\;2^{\mathrm s}$ & $327^\circ\;2'\;13''3$ & $14^\circ\;5'\;43''4$\\
August 16 & $10\;36\;23$ & $319\;50\;54.7$ & $.\;.\;.\;.$\\
19 & $9\;31\;0$ & $319\;17\;10.6$ & $10\;0\;8.4$\\
26 & $9\;23\;7$ & $317\;59\;52.0$ & $8\;44\;35.3$\\
\end{tabular}\end{center}

\Needspace{18\baselineskip}\par\indent Ausserdem wurden dem Prof. \textsc{Gauss} durch Hrn. \textsc{Burckhardt} und Hrn.
\textsc{von Lindenau} noch folgende schätzbare Beobachtungen mitgetheilt:

\begin{center}Meridianbeobachtungen der Pallas auf der Sternwarte der Militärschule in Paris.\end{center}
\begin{center}\fontsize{10}{12}\selectfont\setlength{\tabcolsep}{5pt}\renewcommand{\arraystretch}{1.15}\begin{tabular}{l|c|c|c}
1813 & Sternzeit & Gerade Aufsteig. & Nordl. Abweich.\\\hline
August 4 & $21^{\mathrm u}\;28^{\mathrm m}\;30^{\mathrm s}5$ & $322^\circ\;7'\;38''0$ & $12^\circ\;11'\;58''5$\\
9 & $21\;24\;42$ & $.\;.\;.\;.$ & $11\;32\;38.4$\\
11 & $.\;.\;.\;.$ & $.\;.\;.\;.$ & $11\;15\;28.1$\\
15 & $21\;20\;7.2$ & $320\;1\;48.0$ & $10\;38\;48.2$\\
17 & $21\;18\;35.2$ & $319\;38\;48.5$ & $10\;19\;29.3$\\
19 & $21\;17\;4.3$ & $319\;16\;4.5$ & $9\;59\;18.3$\\
26 & $21\;11\;56.1$ & $317\;59\;1.8$ & $8\;43\;23.4$\\
\end{tabular}\end{center}

\Needspace{11\baselineskip}\begin{center}\leavevmode\SourceLeaf{375}Beobachtungen der Pallas auf der Seeberger Sternwarte:\end{center}

\begin{center}\fontsize{10}{12}\selectfont\setlength{\tabcolsep}{5pt}\renewcommand{\arraystretch}{1.15}\begin{tabular}{l|c|c}
1813 & Mittlere Zeit & Gerade Aufsteig.\\\hline
Sept. 1 & $10^{\mathrm u}\;25^{\mathrm m}\;16^{\mathrm s}6$ & $316^\circ\;58'\;15''9$\\
3 & $10\;16\;9.6$ & $316\;39\;25.5$\\
4 & $10\;11\;35.4$ & $316\;29\;48.7$\\
\end{tabular}\end{center}

\Needspace{25\baselineskip}\par\indent Prof. \textsc{Gauss} übergab diese Beobachtungen dem Hrn. \textsc{Nicolai}. Dieser
geschickte junge Astronom, welcher nunmehr als Gehülfe des Hrn. \textsc{von Lindenau} bei der
Seeberger Sternwarte angestellt ist, verglich sie zuvörderst mit den zuletzt verbesserten
Elementen auf das sorgfältigste, und fand folgende schöne Uebereinstimmung:

\begin{center}Unterschied der Rechnung:\end{center}
\begin{center}\fontsize{10}{12}\selectfont\setlength{\tabcolsep}{7pt}\renewcommand{\arraystretch}{1.15}\begin{tabular}{l|c|c|l}
1813 & \shortstack{Gerade\\Aufsteig.} & \shortstack{Abwei-\\chung} & Beobachter\\\hline
Juni 28 & $+23''6$ & $+22''4$ & Gauss\\
Aug. 4 & $+36.9$ & $-5.9$ & Burckhardt\\
9 & $.\;.\;.\;.$ & $+5.1$ & Burckhardt\\
11 & $.\;.\;.\;.$ & $+6.2$ & Burckhardt\\
15 & $+26.9$ & $+3.7$ & Burckhardt\\
16 & $+35.6$ & $.\;.\;.\;.$ & Gauss\\
17 & $+28.2$ & $-6.6$ & Burckhardt\\
19 & $+28.5$ & $+4.5$ & Gauss\\
19 & $+24.8$ & $-7.4$ & Burckhardt\\
26 & $+22.4$ & $-4.3$ & Gauss\\
26 & $+18.3$ & $+11.1$ & Burckhardt\\
Sept. 1 & $+22.8$ & $.\;.\;.\;.$ & v. Lindenau\\
3 & $+20.9$ & $.\;.\;.\;.$ & v. Lindenau\\
4 & $+49.9$ & $.\;.\;.\;.$ & v. Lindenau\\
\end{tabular}\end{center}

\Needspace{13\baselineskip}\par\indent Aus den Beobachtungen vom 15., 17. und 19. August leitete Hr. \textsc{Nicolai} folgendes
Resultat für die Opposition ab:

\begin{center}\emph{Neunte beobachtete Opposition der Pallas.}\end{center}
\begin{center}1813 August 18. $8^{\mathrm u}\;41^{\mathrm m}\;5^{\mathrm s}$ mittlere Zeit in Göttingen\end{center}
\begin{center}\fontsize{10}{12}\selectfont\setlength{\tabcolsep}{3pt}\renewcommand{\arraystretch}{1.15}\begin{tabular}{lcr}
Wahre Länge der Pallas & $.\;.\;.\;.\;.\;.\;.\;.$ & $325^\circ\;23'\;56''5$\\
Wahre geocentrische Breite & $.\;.\;.\;.\;.\;.\;.$ & $24\;37\;36.1\ \mathrm{N.}$\\
\end{tabular}\end{center}

\noindent Man vergleiche damit die Vorausbestimmung im 127. Stück dieser Anz. vom vorigen Jahre
[S. 356 d. B.].

\par\indent Eine neue Verbesserung der Elemente, um sie auch dieser Opposition noch besser anzupassen,
schien dem Prof. \textsc{Gauss} unter diesen Umständen nicht der Mühe werth zu sein.
